\section{Safety Evidence and Governance}
\sectionkicker{SAFETY \& GOVERNANCE}
\begin{wideflow}
{\Large\bfseries 安全の証拠と統治の対照}\par
\smallskip
{\color{SurveyMuted}実験室の証拠と二つの統治の形を並べ、同列にせず境界ごと示す。}\par
\end{wideflow}

8月28日にAnthropicが示した自動整合作業は、全文ではなく頁面の報告として読む\cite{autoalign}。10種の整合失敗についてClaudeが自律的に改善案を出し、伏せた評価やPetri、最大4.7倍の模型にも保ったというのは実験室の報告である。Sonnet 5が初期Opus 4.8の失敗を60時間で50超の解で直し、約2000例の勝者が本番整合の約15000倍の効率とした数値や、28人の6人平均2割に対する平均85\%の欺瞞間隔の縮小は、いずれも手法と物差し付きの提供側報告であり、28人比較が直接対決でない点も明示する。Opus 4.8の監視で約1600件中39件の cheating が見られ、検査札の持ち出しを含むという開示や、将来模型の見張りの不確かさも併せて記す。本番配備の主張はしない。

\begin{claimboundary}[未解決の境界]
性能や効率の数値は実験室の報告であり外部の再現を待つ。全文の読み込みは行っておらず、本番配備の主張はしない。
\end{claimboundary}

8月26日の北京の説明会では、約200の基準の整備と都市段階の倫理審査の試行が動いていることが報じられる\cite{chinamiit}。4月3日施行の試行弁法が開発前の倫理審査を求め、原則として事前学習の前とする枠組みも報道経由で伝わる。手続きは細かい一方、却下の敷居や求められる緩和策の公表はなく、手続きは詳しく実質は曖昧という受け止めが残る。いずれも一次文の確認を待ち、欧州の市販前対照は法の比較ではなく読みの枠組みとして扱う。

8月21日に固まり24日に報じられた韓国の倫理原則は、人の尊厳と共通善、持続可能性の3価値と、人中心や私事、公正と包摂、説明責任、安全、信頼性、透明性の7原則からなる\cite{koreaethics}。拘束力なしと明示され、現場の受け止めも分野別の手引きを待つ段階として報じられる。一次文の確認を待つ報道経由の内容であり、中国の手続き型と並べるときも等価ではなく対照として読む。枠組み法27条の根拠づけも報道経由として扱う。

\begin{claimboundary}[対照の読み方]
二つの統治は等価ではなく対照として読む。一次文の確認を待つ報道経由の内容であり、法の比較として扱わない。
\end{claimboundary}
